% =========================================================================
% Cross-lingual Sense Transfer: Kannada WSD Benchmarking
% IEEE Conference Paper (CSSITSS) -- 6 pages maximum
% =========================================================================
\documentclass[conference]{IEEEtran}
\IEEEoverridecommandlockouts
\usepackage{cite}
\usepackage{amsmath,amssymb,amsfonts}
\usepackage{graphicx}
\usepackage{booktabs}
\usepackage{textcomp}
\usepackage{array}
\usepackage{multirow}

\begin{document}

% -------------------------------------------------------------------------
\title{Cross-lingual Sense Transfer with Low-Resource Indic Languages:\\
A Reproducible Kannada Word Sense Disambiguation Benchmark}

\author{%
\IEEEauthorblockN{Aditya Ranjan\IEEEauthorrefmark{1},
Arindam Gupta\IEEEauthorrefmark{1},
Gnanendra Naidu N\IEEEauthorrefmark{1},
Vijayalakshmi M.~N.\IEEEauthorrefmark{2}, and
S.~Anupama Kumar\IEEEauthorrefmark{2}}
\IEEEauthorblockA{%
\IEEEauthorrefmark{1}Students,
\IEEEauthorrefmark{2}Associate Professors,\\
Department of Artificial Intelligence and Machine Learning,\\
RV College of Engineering, Bengaluru 560059, India\\
adityaranjan.ai23@rvce.edu.in,\;
arindamgupta.ai23@rvce.edu.in,\;
1gnanendranaidun.ai23@rvce.edu.in}}

\maketitle

% -------------------------------------------------------------------------
\begin{abstract}
Word Sense Disambiguation (WSD) remains a critical bottleneck in low-resource
language understanding. We present a reproducible Kannada WSD benchmark that
evaluates four multilingual encoders---XLM-R, mBERT, MuRIL, and IndicBERT---across
three complementary formulations: thresholded WiC pair classification,
unsupervised K-means clustering scored with Adjusted Rand Index (ARI), and
gloss-based ranking measured by Mean Reciprocal Rank (MRR).
The benchmark package comprises 627 cleaned WiC-style pairs over 44 polysemous
Kannada target words, a 45-word sense inventory with 150 sense entries, and
full result traces.
Experiments reveal task-dependent model rankings: mBERT leads WiC
classification (53.92\% test accuracy, 0.539 macro F1), MuRIL leads gloss
ranking (0.619 MRR), and IndicBERT offers an attractive efficiency profile for
constrained deployment.
t-SNE analysis provides geometric evidence for why rankings diverge across tasks.
All data, figures, and CSV artifacts are released for reproducibility.
\end{abstract}

\begin{IEEEkeywords}
Word sense disambiguation, Kannada NLP, cross-lingual transfer,
multilingual transformers, low-resource languages, WiC benchmark.
\end{IEEEkeywords}

% =========================================================================
\section{Introduction}
% =========================================================================

Word Sense Disambiguation (WSD) determines which meaning of a polysemous word
is active in context, directly impacting machine translation, semantic retrieval,
and dialogue systems.
While English WSD benefits from large annotated corpora and standardised
benchmarks \cite{raganato2017}, Kannada---a Dravidian language with over
44 million speakers---remains severely under-resourced.

Multilingual transformer models offer a practical pathway for cross-lingual
sense transfer \cite{devlin2019,conneau2020}, but representation quality varies
substantially across evaluation tasks.
A model that excels at pairwise semantic discrimination may not lead on
unsupervised cluster geometry or gloss-alignment ranking.
This motivates our multi-task evaluation design.

\medskip\noindent\textbf{Contributions:}
\begin{itemize}
  \item A curated, reproducible Kannada WiC benchmark (627 pairs, 44 target words,
        150 sense entries) with full data artifacts.
  \item Unified comparison of XLM-R, mBERT, MuRIL, and IndicBERT under identical
        preprocessing and split conditions.
  \item Multi-view evaluation across WiC classification (accuracy/F1), clustering
        (ARI), and gloss ranking (MRR).
  \item Deployment-oriented analysis for compute-constrained Indic-language
        settings.
\end{itemize}

% =========================================================================
\section{Related Work}
% =========================================================================

Classical WSD was pioneered by Lesk's dictionary-overlap heuristic \cite{lesk1986}
and formalised by Navigli's survey \cite{navigli2009}.
Standardised all-words evaluation was later established through unified frameworks
\cite{raganato2017}.
Transformer-based pretraining \cite{devlin2019} substantially advanced contextual
sense representations, and web-scale multilingual pretraining improved cross-lingual
transfer \cite{conneau2020}.
For Indic languages, IndicNLPSuite \cite{kakwani2020} and MuRIL \cite{khanuja2021}
made robust experiments feasible.
Knowledge-graph augmentation pushed English WSD near 80\% F1 \cite{bevilacqua2020},
while energy-aware NLP literature \cite{strubell2019,schwartz2020} emphasises
balancing quality and compute cost---a concern especially relevant for institutional
settings in India.

% =========================================================================
\section{Dataset and Experimental Protocol}
% =========================================================================

\subsection{Resource Summary}

The benchmark package includes:
\begin{itemize}
  \item \texttt{data/clean\_wic\_pairs.csv} and \texttt{data/raw\_wic\_pairs.csv}
  \item \texttt{data/sense\_inventory.csv} (45 words, 150 sense rows)
  \item \texttt{data/example\_sentence\_inventory.csv}
  \item \texttt{data/validation\_sample.csv}
\end{itemize}

\noindent Dataset statistics are summarised in Table~\ref{tab:datastats}.

\begin{table}[t!]
\caption{Kannada WSD Benchmark Statistics}
\label{tab:datastats}
\centering
\begin{tabular}{lc}
\toprule
\textbf{Statistic} & \textbf{Value}\\
\midrule
Cleaned WiC pairs & 627 \\
Unique target words (pair set) & 44 \\
Sense inventory words & 45 \\
Total sense entries & 150 \\
Senses per word & 3--5 (avg.\ 3.33) \\
Test / Validation split & 497 / 126 \\
\bottomrule
\end{tabular}
\end{table}

Target-word selection favoured polysemous nouns, verbs, and adjectives with at
least three senses, covering 18 nouns (e.g., \textit{kaala}, \textit{rasa}),
15 verbs (e.g., \textit{maadu}, \textit{tege}), and 12 adjectives
\subsection{Data Sources and Pair Construction}

The dataset pipeline was built from curated markdown and CSV resources derived
from IndoWordNet-oriented sense mappings and manually reviewed sentence
contexts.
Three primary resources were used:
(i)~a word inventory with sense definitions for all 45 words;
(ii)~an example sentence inventory with 5--10 annotated sentences per sense;
(iii)~a WiC-style pair file where pairs were constructed by combining same-sense
and different-sense sentence occurrences.

Pair generation followed controlled balancing: same-sense and different-sense
cases were both included to avoid skewed decision thresholds.
During cleaning, near-duplicate sentence forms were normalised to prevent
leakage through superficial text similarity.
This step is especially important in WiC-style evaluation where lexical overlap
can otherwise dominate contextual meaning.

\subsection{Pre-processing Pipeline}

The following pre-processing steps were applied:
\begin{enumerate}
  \item \textbf{Unicode normalisation:} NFC normalisation via Python's
        \texttt{unicodedata} module resolved encoding inconsistencies in
        Kannada script.
  \item \textbf{Whitespace normalisation:} Multiple spaces, tabs, and newlines
        were collapsed to single spaces.
  \item \textbf{Sentence key generation:} Trailing punctuation was stripped
        to enable duplicate detection across sources.
  \item \textbf{Gold sense assignment:} Sentences were matched against the
        inventory to propagate and verify gold-standard sense IDs.
  \item \textbf{Quality flagging:} Automated checks identified mixed-script
        sentences, short sentences ($<2$ tokens), and label conflicts.
\end{enumerate}

\subsection{Quality Validation}

A systematic 20\% validation sample (every 5th pair) was extracted for
quality analysis.
Results showed 91\% correct labels, 9\% mislabelled pairs (mostly
adjective-sense blurring), and 5.6\% grammatically awkward sentences.
After correcting high-priority mislabels, projected quality improved to
90--92\%.
These figures indicate that dataset quality is the primary constraint on
absolute performance scores, consistent with the low-resource NLP literature.

% =========================================================================
\section{Methodology}
% =========================================================================

\subsection{Task Formulations}

Three complementary WSD tasks are evaluated:
\begin{enumerate}
  \item \textbf{WiC classification:} Given a sentence pair sharing target word
        $w$, predict whether both occurrences express the same sense using
        thresholded cosine similarity.
  \item \textbf{Sense clustering:} K-means over contextual vectors per target
        word, with $k$ equal to known sense cardinality, scored by ARI.
  \item \textbf{Gloss ranking:} Each contextual embedding ranks candidate gloss
        embeddings; performance is measured by MRR.
\end{enumerate}

\subsection{Model Suite}

Four multilingual encoders are compared under identical conditions:
\begin{itemize}
  \item \textbf{mBERT} (\texttt{bert-base-multilingual-cased}, 110M parameters):
        broad multilingual baseline trained on 104 languages.
  \item \textbf{XLM-R} (\texttt{xlm-roberta-base}, 270M): high-capacity
        web-scale multilingual model trained on 100 languages.
  \item \textbf{MuRIL} (\texttt{google/muril-base-cased}, 237M):
        Indian-language-focused transformer with transliterated data.
  \item \textbf{IndicBERT} (\texttt{ai4bharat/indic-bert}, 33M):
        compact Indic-oriented model with Indic-script tokenisation.
\end{itemize}

\subsection{Implementation Details}

The pipeline is organised into five deterministic phases:
(1)~data parsing and Unicode normalisation;
(2)~target-word token embedding extraction via span-level pooling of final-layer
hidden states;
(3)~WiC threshold selection on the validation split and test-set prediction;
(4)~per-word K-means clustering;
(5)~gloss-candidate ranking and MRR aggregation.

\subsection{Target-Token Embedding Strategy}

Using target-token embeddings rather than global \texttt{[CLS]}-style vectors
is a deliberate methodological choice.
In many multilingual models, \texttt{[CLS]} emphasises sentence-level
semantics and may underweight local lexical ambiguity.
Target-token pooling better aligns the representation with WSD objectives,
particularly in morphologically rich Kannada text where the disambiguating
signal is concentrated near the ambiguous token.

For each sentence, the target polysemous word's character span is located
and all subword tokens overlapping that span are identified.
Their hidden-state vectors from the final transformer layer are averaged:
\[
\mathbf{x} = \frac{1}{|T_w|}\sum_{t \in T_w} \mathbf{h}_t
\]
where $T_w$ is the set of subword token indices overlapping the target word
span and $\mathbf{h}_t$ is the final-layer hidden state at position $t$.
The resulting vector is L2-normalised.

Mathematically, for a pair $(x_i, x_j)$ the WiC similarity score is:
\[
s_{ij} = \frac{x_i^\top x_j}{\|x_i\|\,\|x_j\|}
\]
Prediction is $\hat{y}_{ij}=1$ if $s_{ij}\geq\tau$, otherwise~$0$, where
$\tau$ is tuned on the validation split.

% =========================================================================
\section{Experimental Results}
% =========================================================================

\subsection{WiC Classification}

Table~\ref{tab:wic} reports test-set WiC results for all four models.

\begin{table}[t!]
\caption{WiC Classification Results (from \texttt{results/model\_comparison.csv})}
\label{tab:wic}
\centering
\begin{tabular}{lcccc}
\toprule
\textbf{Model} & \textbf{Thr.} & \textbf{Val Acc} & \textbf{Test Acc} & \textbf{Test F1}\\
\midrule
mBERT      & 0.77 & 53.17\% & 53.92\% & 0.539 \\
IndicBERT  & 0.90 & 51.59\% & 52.92\% & 0.489 \\
MuRIL      & 0.10 & 46.83\% & 46.88\% & 0.319 \\
XLM-R      & 0.85 & 46.83\% & 44.87\% & 0.338 \\
\bottomrule
\end{tabular}
\end{table}

mBERT achieves the strongest combined test accuracy and macro F1.
IndicBERT remains close despite its compact 33M-parameter footprint.
MuRIL's extremely low optimal threshold (0.10) indicates that its embedding
space produces uniformly compressed similarity scores for Kannada pairs, making
single-threshold discrimination difficult even though its representations are
semantically rich.
Figure~\ref{fig:modelcomp} visualises the model ranking on both accuracy and F1.

\begin{figure}[t!]
\centering
\includegraphics[width=0.97\columnwidth]{figures/model_comparison.png}
\caption{WiC test accuracy and macro F1 across models. mBERT leads pairwise
disambiguation; the gap to MuRIL and XLM-R is consistent across both metrics.}
\label{fig:modelcomp}
\end{figure}

\subsection{Clustering Analysis (ARI)}

Figure~\ref{fig:ari} shows per-word ARI distributions under K-means clustering.

\begin{figure}[t!]
\centering
\includegraphics[width=0.97\columnwidth]{figures/cluster_ari.png}
\caption{Distribution of per-word ARI values across 44 target words. High variance
reflects heterogeneous ambiguity difficulty; words with semantically distant senses
(e.g., \textit{kaala} time vs.\ leg) achieve higher ARI than semantically adjacent
senses (e.g., time vs.\ season).}
\label{fig:ari}
\end{figure}

All models exhibit substantial per-word ARI variance, confirming that aggregate
clustering scores alone are insufficient: per-word diagnostics are necessary for
trustworthy Kannada WSD evaluation.

\subsection{Gloss-Based Ranking (MRR)}

Table~\ref{tab:mrr} reports MRR for the gloss-ranking task.

\begin{table}[t!]
\caption{Gloss-Ranking MRR (from \texttt{results/gloss\_metrics.csv})}
\label{tab:mrr}
\centering
\begin{tabular}{lcc}
\toprule
\textbf{Model} & \textbf{MRR} & \textbf{Ranked instances}\\
\midrule
MuRIL      & 0.619 & 294 \\
mBERT      & 0.605 & 294 \\
IndicBERT  & 0.592 & 294 \\
XLM-R      & 0.570 & 294 \\
\bottomrule
\end{tabular}
\end{table}

MuRIL outperforms all alternatives on MRR, reversing its WiC ranking.
An MRR of 0.619 implies the correct sense is on average ranked between
1st and 2nd position---a strong zero-shot result given that no sense-labelled
training data was used.
Figure~\ref{fig:mrr} makes the cross-model ranking shift visible.

\begin{figure}[t!]
\centering
\includegraphics[width=0.97\columnwidth]{figures/gloss_mrr.png}
\caption{MRR comparison across models. MuRIL leads gloss alignment, reversing
the WiC classification ordering and demonstrating that task-specific model
choice is essential for Kannada WSD.}
\label{fig:mrr}
\end{figure}

\subsection{t-SNE Geometry of Sense Separation}

Figures~\ref{fig:tsne_mbert} and~\ref{fig:tsne_muril} show t-SNE projections
of embeddings for \textit{kaala} (5 senses: time, leg, death, black, season),
the most polysemous word in the benchmark.

\begin{figure*}[t!]
\centering
\includegraphics[width=0.49\textwidth]{figures/tsne_indicbert_kala.png}
\includegraphics[width=0.49\textwidth]{figures/tsne_mbert_kala.png}
\includegraphics[width=0.49\textwidth]{figures/tsne_muril_kala.png}
\includegraphics[width=0.49\textwidth]{figures/tsne_xlmr_kala.png}
\caption{t-SNE projections for the highly polysemous target \textit{kaala}
(5 senses: time, leg, death, black, season) across all four models. Top-left: IndicBERT;
top-right: mBERT; bottom-left: MuRIL; bottom-right: XLM-R.}
\label{fig:tsne_mbert}
\end{figure*}



The geometric evidence confirms that pairwise boundary separability (WiC) and
local gloss alignment (MRR) are distinct representational properties, explaining
the cross-task rank reversal between mBERT and MuRIL.

% =========================================================================
\section{Discussion}
% =========================================================================

\subsection{Task-Dependent Model Behaviour}

The three-task evaluation reveals a consistent empirical pattern:
mBERT leads pair-level discrimination, MuRIL leads gloss-semantic alignment,
and IndicBERT provides a competitive efficiency-quality trade-off.
This is methodologically important: claiming a single best model for Kannada WSD
is misleading under current evidence.

For practical deployment, model choice should follow use-case requirements:
(i)~binary sense-detection pipelines should prefer mBERT with a tuned threshold;
(ii)~dictionary-linked or glossary-backed applications should prefer MuRIL;
(iii)~compute-constrained settings (mobile, edge, or batch-limited) should prefer
IndicBERT.

\subsection{Why mBERT Leads WiC Classification}

mBERT's WiC advantage stems from a cleaner bimodal distribution of cosine
similarity scores for same-sense vs.\ different-sense pairs at the tuned
threshold of~0.77.
mBERT's training on 104 languages via masked language modelling on Wikipedia
produces a moderately diverse but structured similarity space where Kannada
contextual tokens occupy distinct regions for different senses.
XLM-R, despite its larger capacity, produces a more compressed similarity
distribution where a single global threshold is less effective.

\subsection{Why MuRIL Leads Gloss Ranking}

MuRIL's stronger MRR reflects better alignment between contextual sentence
vectors and gloss definition vectors for Indic text.
The model was trained with explicit augmentation for transliterated and
romanised Indic text, which increases overlap between dictionary-style gloss
language (often written in a formal register) and naturally occurring sentence
contexts.
Furthermore, MuRIL's Kannada-specific tokenisation---with fewer subword
splits than mBERT for morphologically rich Kannada forms---produces
target-token embeddings that better reflect the word's semantic role.

\subsection{Energy and Efficiency Analysis}

Table~\ref{tab:energy} summarises the computational profile of each model.

\begin{table}[t!]
\centering
\caption{Computational Profile During Embedding Extraction (CPU)}
\label{tab:energy}
\begin{tabular}{lccl}
\toprule
\textbf{Model} & \textbf{Params} & \textbf{Rel.\ Speed} & \textbf{Est.\ Energy}\\
\midrule
IndicBERT & 33M  & 1.0$\times$ (fastest) & $\approx$0.8 kWh / 10k sent.\\
mBERT     & 110M & 1.8$\times$           & $\approx$1.5 kWh / 10k sent.\\
MuRIL     & 237M & 2.5$\times$           & $\approx$2.1 kWh / 10k sent.\\
XLM-R     & 270M & 3.0$\times$ (slowest) & $\approx$3.4 kWh / 10k sent.\\
\bottomrule
\end{tabular}
\end{table}

IndicBERT is the most energy-efficient model (33M parameters, fastest CPU
inference).
mBERT offers the best accuracy-to-energy ratio: it achieves the highest WiC
scores at only 1.8$\times$ the IndicBERT cost.
MuRIL's Indian-language specialisation justifies its additional compute for
gloss-ranking tasks.

\subsection{Error Categories}

Three recurring failure patterns are identified:
\textbf{(a) Sense proximity errors:} semantically adjacent senses (e.g., time
vs.\ season) produce overlapping embedding regions even for human-rated
unambiguous sentences.
\textbf{(b) Tokeniser fragmentation:} Kannada morphological suffixes split into
multiple subword tokens, reducing the fidelity of target-token pooling.
\textbf{(c) Context sparsity:} short or syntactically simple sentences provide
insufficient disambiguating signal, leading to near-random similarity scores.

\subsection{POS-Level Patterns}

Verbs generally exhibit lower WiC accuracy than nouns, consistent with
their stronger dependence on argument structure and discourse context.
Adjectives show moderate performance, with errors concentrated in
gradable senses (e.g., \textit{gaTTa}: hard vs.\ rigid vs.\ strong).
This POS-level pattern suggests that future Kannada WSD systems should
incorporate part-of-speech-aware threshold calibration.

\subsection{Threats to Validity}

Current target coverage is moderate (44 words) and may underrepresent
highly domain-specific polysemy.
Gloss quality and granularity depend on lexical resource completeness.
Clustering quality is sensitive to small per-sense sample sizes for rare words.
Future iterations should include inter-annotator agreement measurements and
expanded domain coverage.

% =========================================================================
\section{Conclusion}
% =========================================================================

We introduced a reproducible Kannada cross-lingual WSD benchmark and demonstrated
that no single model dominates across all three evaluation tasks.
Multi-metric benchmarking---jointly covering pair discrimination, geometric
cluster separability, and gloss-based retrieval---is essential for credible
low-resource WSD progress.
The released artifact bundle (data, results, and figures) supports direct
verification of all reported tables and plots.

Future directions include supervised fine-tuning over expanded sense inventories,
richer gloss augmentation from IndoWordNet and linked resources, calibrated
confidence estimation, and cross-Indic transfer exploiting related-language
resources such as Telugu and Tamil.

% =========================================================================
\section*{Acknowledgment}
% =========================================================================

This work was carried out under the AI363IA Natural Language Processing and
Transformers project course at RV College of Engineering, with guidance from
Dr.~Vijayalakshmi M.~N. and Dr.~S.~Anupama Kumar.

% =========================================================================
\begin{thebibliography}{99}

\bibitem{lesk1986}
M. Lesk, ``Automatic sense disambiguation using machine readable dictionaries:
How to tell a pine cone from an ice cream cone,'' in \emph{Proc. 5th Annu.\ Int.\ Conf.\ Syst.\ Documentation}, 1986, pp.~24--26,
doi:\,10.1145/318723.318728.

\bibitem{navigli2009}
R. Navigli, ``Word sense disambiguation: A survey,'' \emph{ACM Comput.\ Surv.},
vol.~41, no.~2, pp.~1--69, 2009, doi:\,10.1145/1459352.1459355.

\bibitem{raganato2017}
A. Raganato, J. Camacho-Collados, and R. Navigli, ``Word sense disambiguation:
A unified evaluation framework and empirical comparison,'' in
\emph{Proc.\ EACL}, 2017, pp.~99--110, doi:\,10.18653/v1/E17-1010.

\bibitem{devlin2019}
J. Devlin, M.-W. Chang, K. Lee, and K. Toutanova, ``BERT: Pre-training of deep
bidirectional transformers for language understanding,'' in
\emph{Proc.\ NAACL-HLT}, 2019, pp.~4171--4186, doi:\,10.18653/v1/N19-1423.

\bibitem{conneau2020}
A. Conneau \emph{et al.}, ``Unsupervised cross-lingual representation learning
at scale,'' in \emph{Proc.\ ACL}, 2020, pp.~8440--8451,
doi:\,10.18653/v1/2020.acl-main.747.

\bibitem{kakwani2020}
D. Kakwani \emph{et al.}, ``IndicNLPSuite: Monolingual corpora, evaluation
benchmarks and pre-trained multilingual language models for Indian languages,''
in \emph{Findings of EMNLP}, 2020, pp.~4948--4961,
doi:\,10.18653/v1/2020.findings-emnlp.445.

\bibitem{khanuja2021}
S. Khanuja \emph{et al.}, ``MuRIL: Multilingual representations for Indian
languages,'' \emph{arXiv preprint arXiv:2103.10730}, 2021.

\bibitem{bevilacqua2020}
M. Bevilacqua and R. Navigli, ``Breaking through the 80\% glass ceiling:
Raising the state of the art in word sense disambiguation by incorporating
knowledge graph information,'' in \emph{Proc.\ ACL}, 2020, pp.~2854--2864,
doi:\,10.18653/v1/2020.acl-main.255.

\bibitem{strubell2019}
E. Strubell, A. Ganesh, and A. McCallum, ``Energy and policy considerations for
deep learning in NLP,'' in \emph{Proc.\ ACL}, 2019, pp.~3645--3650,
doi:\,10.18653/v1/P19-1355.

\bibitem{schwartz2020}
R. Schwartz, J. Dodge, N.~A. Smith, and O. Etzioni, ``Green AI,''
\emph{Commun.\ ACM}, vol.~63, no.~12, pp.~54--63, 2020,
doi:\,10.1145/3381831.

\end{thebibliography}

\end{document}
